\documentclass[11pt]{article}
\usepackage[a4paper,margin=14mm]{geometry}
\usepackage[no-math]{fontspec}
\setmainfont{Latin Modern Roman}
\usepackage{amsmath,amssymb,amsthm,xfrac,xspace,hyperref,graphicx,comment}
\excludecomment{DefTralics}

\providecommand{\bibliofont}{\small}
\newcommand{\ie}{i.e.,\xspace}
\newcommand{\eg}{e.g.,\xspace}
\newcommand{\etc}{etc.\xspace}
\newcommand{\cf}{cf.\xspace}
\newcommand{\resp}{resp.\xspace}
\newcommand{\smaller}{\small}
\newcommand{\Subsection}{\subsection}
\newcommand{\Subsubsection}{\subsubsection}
\newcommand{\skpt}{\smallskip}
\emergencystretch=3em
\numberwithin{equation}{section}\input{author-preamble.tex}
\begin{document}
\begin{center}{\Large Stable item 2047}\end{center}
\paragraph{Source and attribution.}
Antoine Henrot, Ilaria Lucardesi, \emph{Body of constant width with minimal area in a given annulus}, Journal de l’École polytechnique — Mathématiques 8 (2021), publisher version, DOI 10.5802/jep.150. \href{https://jep.centre-mersenne.org/articles/10.5802/jep.150/}{Publisher article}. Authors' text: \href{https://creativecommons.org/licenses/by/4.0/}{CC BY 4.0}.
\paragraph{Editorial problem boundary.}
Prove only Theorem 4.1 below in the class of Reuleaux polygons of unit width, prescribed inradius \(1-1/\sqrt3\le r<1/2\), and at most \(2N+1\) sides. The unique minimiser is the exact polygon \(\Omega(r)\), up to rigid motion, when the class is nonempty. The full shape-derivative, rigid-cluster and analytic concavity construction is retained. No extension to arbitrary constant-width bodies through the separate density argument is selected.
\paragraph{Difficulty (editorial): H2.}
A fixed-inradius constraint prevents the usual unconstrained Blaschke deformation argument from identifying the area minimiser directly. The author proves rigidity of the admissible clusters, derives their exact area formula and strict interior concavity, and uses the constrained arc-count optimisation to exclude competing polygon configurations. This geometric reduction and analytic sign construction establish the unique polygon optimiser without relying on the separate density claim for arbitrary constant-width bodies.
\paragraph{Editorial transcription note.}
Original author language, mathematics and ordering are preserved. Publisher front matter is replaced by this independent article wrapper. Bibliography is recovered from matching cached publisher HTML, including its TeX formulas. No new proof or mathematical repair is supplied. Surrounding original results are retained for conservative context and are not extra selected corpus claims.
\clearpage
\input{author-body.tex}
\begin{thebibliography}{99}
\bibitem{Be} A. S. Besicovitch - “Minimum area of a set of constant width” , in Proc. Sympos. Pure Math., Vol. VII , American Mathematical Society, Providence, RI, 1963, p. 13-14
\bibitem{Bla} W. Blaschke - “Konvexe Bereiche gegebener konstanter Breite und kleinsten Inhalts” , Math. Ann. 76 (1915) no. 4, p. 504-513
\bibitem{BF} T. Bonnesen \& W. Fenchel - Theorie der konvexen Körper , Berichtigter Reprint , Springer-Verlag, Berlin-New York, 1974, English transl.: Theory of convex bodies , BCS Associates, Moscow, ID, 1987
\bibitem{Buc} H. Bückner - “Über Flächen von fester Breite” , Jahresber. Deutsch. Math.-Verein. 46 (1936), p. 96-139
\bibitem{CCG} S. Campi, A. Colesanti \& P. Gronchi - “Minimum problems for volumes of convex bodies” , in Partial differential equations and applications , Lecture Notes in Pure and Appl. Math. , vol. 177 , Dekker, New York, 1996, p. 43-55
\bibitem{CG} G. D. Chakerian \& H. Groemer - “Convex bodies of constant width” , in Convexity and its applications , Birkhäuser, Basel, 1983, p. 49-96
\bibitem{Eg} H. G. Eggleston - “A proof of Blaschke’s theorem on the Reuleaux triangle” , Quart. J. Math. Oxford Ser. (2) 3 (1952), p. 296-297
\bibitem{Ga} M. Ghandehari - “An optimal control formulation of the Blaschke-Lebesgue theorem” , J. Math. Anal. Appl. 200 (1996) no. 2, p. 322-331
\bibitem{Ha} E. M. Harrell - “A direct proof of a theorem of Blaschke and Lebesgue” , J. Geom. Anal. 12 (2002) no. 1, p. 81-88
\bibitem{H} A. Henrot - Extremum problems for eigenvalues of elliptic operators , Frontiers in Mathematics , Birkhäuser Verlag, Basel, 2006
\bibitem{HLu} A. Henrot \& I. Lucardesi - “A Blaschke-Lebesgue theorem for the Cheeger constant” , 2020
\bibitem{HP} A. Henrot \& M. Pierre - Variation et optimisation de formes. Une analyse géométrique , Mathématiques \& Applications , vol. 48 , Springer, Berlin, 2005
\bibitem{KM} Y. S. Kupitz \& H. Martini - “On the isoperimetric inequalities for Reuleaux polygons” , J. Geom. 68 (2000) no. 1-2, p. 171-191
\bibitem{Le} H. Lebesgue - “Sur le problème des isopérimètres et sur les domaines de largeur constante” , Bull. Soc. Math. France C.R. 7 (1914), p. 72-76
\bibitem{MMO} H. Martini, L. Montejano \& D. Oliveros - Bodies of constant width. An introduction to convex geometry with applications , Birkhäuser/Springer, Cham, 2019
\bibitem{Ma2} A. E. Mayer - “Über Gleichdicke kleinsten Flächeninhalts” , Anzeiger Wien 71 (1934), p. 4
\bibitem{Ma1} A. E. Mayer - “Der Inhalt der Gleichdicke” , Math. Ann. 110 (1935) no. 1, p. 97-127
\end{thebibliography}
\end{document}
